\documentclass[10pt,a4paper]{amsart}
\usepackage[margin=19mm]{geometry}
\usepackage{amsmath,amssymb,mathtools}
\usepackage[scr=rsfs,cal=euler]{mathalfa}
\usepackage{xfrac}
\usepackage[unicode,hidelinks,bookmarks=false]{hyperref}
\newcommand{\ie}{i.e.\ }
\newcommand{\cf}{cf.\ }
\newcommand{\resp}{respectively\ }
\newcommand{\Subsection}{\subsection}
\providecommand{\nobreakdash}{\nobreak\hbox{-}}
\setlength{\emergencystretch}{2em}
\multlinegap0pt
\numberwithin{equation}{section}
\newtheorem{theorem}{Theorem}[section]
\newtheorem{lemma}[theorem]{Lemma}
\newtheorem{proposition}[theorem]{Proposition}
\newtheorem{remark}[theorem]{Remark}
\allowdisplaybreaks
\title[Global existence and uniqueness of strong solutions for a]{Global existence and uniqueness of strong solutions for a viscoelastic coupled system of two wave equations}
\author{Houda Fadel, Djamila Benterki, Salim A. Messaoudi}
\date{Source: Electronic Journal of Differential Equations, Vol. 2026 (2026), No. 16, pp. 1--23; DOI 10.58997/ejde.2026.16; published February 19, 2026; publisher version}
\begin{document}
\maketitle
\noindent\emph{Electronic Journal of Differential Equations}, Vol. 2026 (2026), No. 16, pp. 1--23.\newline
ISSN: 1072-6691. DOI \href{https://doi.org/10.58997/ejde.2026.16}{10.58997/ejde.2026.16}. This work is licensed under a CC BY 4.0 license.\par\smallskip
\noindent\textit{Editorial scope.} Counted claim: original Theorem 3.1 of the published article (native label thm3, source lines 224--234, printed as Theorem 3.1 in the published PDF), namely that under (A1)--(A3) the viscoelastic coupled system has a unique strong solution with the stated regularity for every finite time horizon, together with its complete original author proof at source lines 236--632. Necessary complete context is retained uncounted: the whole Introduction with the coupled system and hypotheses (A1)--(A3) and the whole of Section 2 with Lemmas 2.1 and 2.2 (source lines 73--223). Nothing inside the delivered ranges is omitted, so every printed section, theorem, lemma and equation number of the delivered unit coincides with the published PDF. The delivered unit ends after the counted proof with the reference list of the cited published entries; the figures of the numerical section and every line of the numerical and examples sections lie outside the delivered ranges, so no illustration line had to be dropped anywhere in this unit and no bounded source comparison was needed. Corpus relationship to problem 3413, stated explicitly because both units come from the same article: the published master batch-1140-1165-r002 already contains problem\_id 3413 with source\_id doi:10.58997/ejde.2026.16 and counted claim 10.58997/ejde.2026.16/Theorem4.10, whose primary statement is source lines 1175--1186 with its proof at source lines 1188--1431; the present unit 7016 counts a different published theorem, Theorem 3.1 with its own complete proof at source lines 236--632, which problem 3413 used only as an uncounted context block (its declared context source lines 224--633) and never as its counted claim. The mathematical subject is different as well: Theorem 3.1 is the global existence and uniqueness of strong solutions for the coupled system built by Faedo-Galerkin approximation, whereas Theorem 4.10 is the general energy decay estimate with an inverse-G2 rate under the convexity assumption on the relaxation function. The two units therefore count disjoint claims of one source, as the corpus precedent for several units per source permits, and they carry different claim keys. Printed slips kept literal: the proof prints the uniqueness argument through the admissible range of the auxiliary parameter exactly as published, and Lemma 4.1 of the article refers to Theorem 3 by a short label, a slip that lies outside the delivered ranges. No mathematics was written, repaired or re-derived.
 \section{Introduction}
 
 Wave equations with memory effects have been a central topic in the analysis of partial
 differential equations, owing to their applications in physics, engineering, and material
 science. The focus has been on the existence, stability, and energy decay of solutions, with significant progress made for both single and coupled wave equations.
 To motivate our work, data collected and results are reviewed to provide a thorough study of
 both the single and coupled viscoelastic wave equation. We start with the pioneer works of Dafermos \cite{Dafermos1, Dafermos2}, where the author showed that solutions decay to zero
 for smooth, monotone relaxation functions $g$, but the decay rate remained unspecified. Cavalcanti et al.\ \cite{CavalcantiDS} studied
\[
|u_t|^\rho u_{tt} - \Delta u - \Delta u_{tt} + \int_0^t g(t - s) \Delta u(s) \, ds - \gamma \Delta u_t = 0,\quad \text{in } \Omega \times \mathbb{R}^{*}_{+}
\]
proving global existence for $\gamma \geq 0$ and exponential energy decay for $\gamma > 0$. Messaoudi and Tatar \cite{MTatar1, MTatar2} extended these results to include source terms, showing exponential and polynomial decay depending on the rate of decay of g(t).
Cavalcanti et al.\ \cite{CavalcantiFerreira} considered a viscoelastic wave equation with a
local frictional damping, where the kernel $g$ satisfies, for two positive constants $\xi_1$
and $\xi_2$,
\begin{align*}
-\xi_1 g(t) \leq g'(t) \leq -\xi_2 g(t), \quad \forall t \in \mathbb{R}_{+}.
\end{align*}
Under the above assumptions and subject to certain restrictions on the control zone, they established an exponential decay result. Berrimi and Messaoudi \cite{Berrimi2004} later
improved this result by proving that the viscoelastic dissipation alone is sufficient to stabilize the system, and subsequently extended it to the case where a source term competes with the dissipation \cite{Berrimi2006}.
Further contributions including the analysis of nonlinear damping with polynomially decaying kernels \cite{Cavalcanti2003} and global existence with uniform decay for source-type problems \cite{Li2011}  appeared later.
 Without restricting $g(t)$ to exponential or polynomial decay, Messaoudi \cite{SMessaoudi1, SMessaoudi2} generalized all the existing decay results by considering the following equation
\[
u_{tt} - \Delta u + \int_0^t g(t - \tau) \Delta u(\tau) \, d\tau = b |u|^{\gamma} u,\quad b \geq 0,
\]
where $g'(t) \leq -\xi(t) g(t)$, for a non-increasing differentiable $\xi(t)$.
These latter results inspired other authors and more results have been established, see \cite{SaidH2011, Hao, PJS, WSH, MessaoudiMustafa}.

Further improvement was establish by Messaoudi and Al-Khulaifi \cite{MessaoudiWaled}, who addressed the problem in \cite{CavalcantiDS}, for $\gamma =0$ and for a kernel satisfying, for $1\leq p\leq \frac{3}{2}$, the condition
\[
g'(t) \leq -\xi(t) g^p(t), \quad \forall t \in \mathbb{R}_{+}.
\]
They proved that the energy of solutions converges to zero at the same rate as $g(t)$, provided $g(t)$ decays faster than $s^{-2}$ at infinity. Later, Mustafa \cite{MM181} extended his earlier joint work with Messaoudi \cite{MM}, which dealt with kernels satisfying $g'(t) \leq -G(g(t))$, to the more general condition $$g'(t) \leq -\xi(t) G(g(t)), \quad \forall t \in \mathbb{R}_{+},$$ where $G$ is strictly increasing and convex function. This extension provided explicit energy decay estimates and encompassed a wider range of decay behaviors.

Coupled wave systems also have consistently drawn attention due to their fundamental role in understanding stability and long-term dynamics. In a recent contribution, Guesmia, Messaoudi, and Zahri \cite{GASM} analyzed an abstract coupled system involving a viscoelastic component interacting with a purely elastic one. Precisely, they looked into the problem
\begin{gather*}
u_{tt}(t) + Au(t) - \int_0^{t} g(s) A^{\theta} u(t - s) \, ds + \alpha u + \beta B v(t) = f(u), \quad t > 0, \\
v_{tt}(t) + A v(t) + \beta B u(t) = 0, \quad t > 0,\\
u(0) = u_0, \quad u_t(0) = u_1, \quad v(0) = v_0, \quad v_t(0) = v_1,
\end{gather*}
 where $A$ is a self-adjoint positive operator on a Hilbert space $H$, $\theta \in [0,1]$,
 $\alpha > 0$, and $g$ is a positive nonincreasing relaxation function.
 Their analysis established general decay rates for the system energy, showing that the
 asymptotic behavior depends strongly on the relaxation function and the coupling parameter.
 Numerical experiments were also provided, highlighting the impact of memory effects on
 the qualitative behavior of solutions. Earlier, Feng and Li \cite{FengBaoweiHaiyan} considered
 a viscoelastic Lam\'e system with strong damping terms. For two distinct relaxation kernels
 $g_1$ and $g_2$, they assumed conditions of the form $g_i'(t) \leq -\xi_i(t) H_i(g_i(t))$
 ($i=1,2$), which led to explicit and general decay results. Their work improved earlier results
 of Beniani, Taouaf, and Benaissa \cite{Beniani}, who had previously established a general
 decay property for similar coupled viscoelastic structures.
Further investigations were carried out by Jin et al. \cite{JKJT}, who studied a coupled system of two abstract evolution equations with finite memory, given by
\begin{gather*}
 u_{tt}(t) + Au(t) - \int_0^{t} g(s) A u(t - s) \, ds + \alpha u + \beta B v(t)
 = f(u), \quad t > 0, \\
v_{tt}(t) + A v(t) + \beta B u(t) = 0, \quad t > 0,
 \end{gather*}
 where $A$ and $B$ are operators, $\alpha \geq 0$, $\beta > 0$, and $f$ is a nonlinear source.
 Under appropriate assumptions on $g$ and structural conditions on $A$, $B$, and $f$,
 they proved polynomial stability with decay rates of order $t^{-1}$. Several extensions and
related contributions in this direction may be found in \cite{ ORHC, AkilM}.
In addition to viscoelastic systems, purely elastic coupled wave models have also been studied.
For instance, a weakly coupled system
\begin{gather*}
u_{tt} - \Delta u + u_t + \alpha v = 0, \quad \text{in } \Omega \times (0, \infty), \\
v_{tt} - \Delta v + \alpha u = 0, \quad \text{in } \Omega \times (0, \infty),
\end{gather*}
was analyzed in \cite{LRF}, where $\Omega \subset \mathbb{R}^n$ is a bounded domain and
$\alpha$ a positive constant. The authors showed that the energy decays polynomially,
with the rate being optimal, and provided numerical validation in the one-dimensional case.
Similar results on stability and decay rates had earlier been obtained for systems with
boundary damping, notably by Komornik et al. \cite{KVBR} and Aassila \cite{Assila1, Assila2},
where exponential stability was achieved in the case of linear damping and polynomial stability
in the case of nonlinear damping mechanisms.

Motivated by the findings in \cite{GASM}, the present work investigates the coupled viscoelastic wave system
\begin{equation} \label{e1.1}
\begin{gathered}
u_{tt} - \Delta u + \int_0^t g(t - s) \Delta u \, ds + \alpha v = 0, \quad
 \text{in } \Omega \times (0, T), \\
v_{tt} - \Delta v + \alpha u = 0, \quad \text{in } \Omega \times (0, T), \\
u = v = 0, \quad \text{on } \partial \Omega \times (0, T), \\
u(0) = u_0, \, u_t(0) = u_1, \; v(0) = v_0, \, v_t(0) = v_1, \quad \text{in } \Omega,
\end{gathered}
\end{equation}
where \( \Omega \) is a bounded domain in \( \mathbb{R}^N, \) \( \alpha > 0 \), and the initial
data belong to suitable spaces. The relaxation function \( g \) satisfies
\( g'(t) \leq -\xi(t) G(g(t)) \), where $G$ is an increasing and convex function near the origin and
$\xi$ is a nonincreasing function. We establish, in details, an existence and uniqueness result,
then we prove a general and explicit result for energy decay and provide some numerical
illustration to validate our theoretical findings.

The structure of this paper is as follows: the next section presents some assumptions,
remarks and functionals needed in our results. In Section 3, we establish the existence of
solutions by combining energy estimates with the Faedo-Galerkin method.
Section 4 presents several technical lemmas and contains the statement and proof of the
decay results.
In section 5, two numerical tests are conducted aiming to verify our decay result.

 \section{Preliminaries}

 This section presents the foundation for our main result. We set
 \begin{gather*}
 (h \star \varphi)(t) = \int_0^t h(t-s) \varphi(s)ds,\
 (h \diamond \varphi)(t) = \int_0^t h(t-s)\| \varphi(t)- \varphi(s)\| ds, \\
 (h \circ \varphi)(t) = \int_0^t h(t-s)\int_\Omega\| \varphi(t)- \varphi(s)\|^2 dx \, ds.
 \end{gather*}
 For the relaxation function $g$, we assume:
 \begin{itemize}
\item[(A1)] $g: \mathbb{R}^+ \to \mathbb{R}^+ $ is a $ C^1$ function satisfying\\
\begin{equation*}
 g(0) > 0, \quad 1-\int_0^{\infty} g(s) \, ds = l > 0. \label{e2.1}
\end{equation*}

\item[(A2)] There exists a $C^1$ function $G :\mathbb{R}^+ \to \mathbb{R}^+ $
which is linear or it is strictly increasing and strictly convex $C^2$ function on
$(0, r]$, $r \leq g(0)$, with $G(0) = G'(0) = 0$, such that
\begin{equation} \label{e2.2}
g'(t) \leq -\xi(t) G(g(t)),\quad \forall t \geq 0,
\end{equation}
where $\xi$ is a positive non-increasing differentiable function.

\item[(A3)] The constant $\alpha $ is such that
\begin{equation} \label{e2.3}
0 < \alpha < \frac{\sqrt{l} }{C_p},
\end{equation}
where $C_p$ is the Poincar\'e constant.
\end{itemize}

\begin{lemma}[\cite{Munoz}] \label{lem1}
For any function $h \in C^1 (\mathbb{R})$ and any $k \in H^1 (0,T)$, we have
$$
(h \star k)(t) k_t
 = -\frac{1}{2} h(t)  \| k\|^2 + \frac{1}{2} (h' \circ k)(t)
 -\frac{1}{2} \frac{d}{dt} \Big\{( h \circ k)(t)
 - \Big( \int_0^t h(\tau) d\tau \Big) \|k\|^2 \Big\}.
$$
\end{lemma}

\begin{lemma}[\cite{SMessaoudi1}] \label{lem2}
 For $u \in H^1_0(\Omega)$, we have
\begin{equation}
\int_\Omega \Big( \int_0^t g(t - s)(u(t) - u(s)) \, ds \Big)^2 dx
\leq (1 - l) C_p^2 (g \circ \nabla u)(t),\label{e2.4}
\end{equation}
where $l$ is given in {\rm (A1)}.
\end{lemma}

\section{Existence of weak solutions}
Now, we examine the well-posedness of problem \eqref{e1.1}. To establish this, we will combine some energy estimates and the Faedo-Galerkin method.


\begin{theorem} \label{thm3}
 Assume that {\rm (A1)--(A3)} hold.
Then, for each $ (u_0,v_0) \in \left( H^2(\Omega) \cap H^1_0(\Omega)\right)^2$ and
$(u_1,v_1) \in \left( H^1_0(\Omega)\right)^2$, problem \eqref{e1.1} has a unique strong solution
$(u, v)$ satisfying, for every $T > 0$,
\begin{gather*}
u, v \in L^{\infty}\left( [0, T), H^2(\Omega) \cap H^1_0(\Omega)\right),\\
u_{t}, v_{t} \in L^2\left( (0, T), H^1_0(\Omega)\right)\cap H^1_0 \left( (0, T)\times \Omega\right),\\
u_{tt}, v_{tt} \in L^{\infty}\left( (0, T), L^2(\Omega)\right).
\end{gather*}
 \end{theorem}

\begin{proof} \quad

\noindent\textbf{Step 1.} Faedo-Galerkin approximation.
We take
\(\{ \phi_j \}_{j=1}^\infty\)
to be the eigenfunctions of the Laplacian operator subject to Dirichlet boundary conditions, satisfying $-\Delta \phi_i = \lambda_i \phi_i $. Then
$\{\phi_j\}_{j=1}^\infty$ forms an orthogonal basis of $ H_0^1(\Omega)$,
\( H^2(\Omega) \cap H_0^1(\Omega) \), and \( L^2(\Omega) \).
Let $V_{N}$ be spanned by the $N$ first functions $ \phi_1, \phi_2, \dots, \phi_N $.
We seek  an approximate solution $(u,v)$ of the form
\[
u^N(x,t) = \sum_{j=1}^N a_j(t) \phi_j(x), \quad
v^N(x,t) = \sum_{j=1}^N b_j(t) \phi_j(x), \quad t\in (0,T),
\]
which satisfies
\begin{gather*}
\int_\Omega u_{tt}^N \phi_j dx + \int_\Omega \nabla u^N .\nabla \phi_j dx - \int_\Omega \int_0^t g(t-s) \nabla u^N(s) .\nabla \phi_j\, ds\,dx+ \alpha \int_\Omega v^N \phi_j \, dx = 0,\\
\int_\Omega v_{tt}^N \phi_j dx + \int_\Omega \nabla v^N .\nabla \phi_j dx+ \alpha \int_\Omega u^N \phi_j \, dx = 0,\quad 1\leq j\leq N
\label{e3.1}
\end{gather*}
and
\begin{equation} \label{e3.2}
\begin{gathered}
 u^{N}(0) = u^{N}_0= \sum_{i=1}^{N} \langle u_0, \phi_i \rangle \phi_i
 \to u_0, \\
  v^{N}(0) = v^{N}_0 = \sum_{i=1}^{N} \langle v_0, \phi_i \rangle \phi_i
 \to v_0 \quad \text{in } H^2(\Omega) \cap H_0^1(\Omega), \\
 u_{t}^N(0) = u_1^N = \sum_{i=1}^{N} \langle u_1, \phi_i \rangle \phi_i \to u_{1}, \\
 v_{t}^N(0) = v_1^N = \sum_{i=1}^{N} \langle v_1, \phi_i \rangle \phi_i \to v_{1} \quad \text{in } H^1_0(\Omega).
 \end{gathered}
\end{equation}
Standard results of ordinary differential equations guarantee the existence of
a unique local solution
$ (u^N, v^N )$ of \eqref{e3.1}-\eqref{e3.2} defined on
$ [0, t_{N}), 0 < t_{N} \leq T $.
Our next step is to obtain a priori estimates that enable us to extend the
solution to the entire interval.
\smallskip

\noindent\textbf{Step 2.} A priori estimate 1.
By multiplying the first equation in \eqref{e3.1} by \(a'_j(t)\) and summing the results over \(j\)
from 1 to \(N\), we obtain, for all \(0 < t \leq t_N\),
\begin{equation}
\begin{aligned}
&\frac{d}{dt} \Big( \frac{1}{2} \| u^N_{t}(t) \|^2 + \frac{1}{2} \| \nabla u^N(t) \|^2 \Big)
+ \alpha \int_{\Omega} v^N(t) \, u^N_{t}(t) \, dx
&- \int_0^t g(t-s) \int_{\Omega} \nabla u^N(s) \nabla u^N_{t}(t) \, dx ds =0.
\end{aligned}\label{e3.3}
\end{equation}
Using Lemma \ref{lem1} for the last term of the left-hand side of \eqref{e3.3}
we find that
\begin{equation}
\begin{aligned}
&\frac{d}{2dt} \Big( \| u^N_{t}(t) \|^2
+ \Big[ 1-\int_0^t g(s) ds \Big] \| \nabla u^N(t) \|^2 + (g\circ \nabla u^N)(t) \Big)
+\alpha \int_{\Omega} v^N(t) \, u^N_{t}(t) dx \\
& +\frac{1}{2} g(t)\parallel \nabla u^N(t) \parallel^2 - \frac{1}{2}( g'\circ\nabla u^N) (t) =0.
\end{aligned} \label{e3.4}
 \end{equation}
Multiplying the second equation in \eqref{e3.1} by $b'_j(t)$ and summing with respect to $j$, we obtain
 \begin{equation}
 \frac{d}{2 dt} \Big( \| v^N_{t}(t) \|^2 + \| \nabla v^N (t) \|^2 \Big)
 + \alpha \int_{\Omega} u^N(t) \, v^N_{t}(t) \, dx =0, \quad 0 < t \leq t_N.
 \label{e3.5}
\end{equation}
Then, combining \eqref{e3.4} and \eqref{e3.5} yields
\begin{equation}
\begin{aligned}
 &\frac{d}{2 dt} \Big( \| u^N_{t}(t) \|^2 + \Big[ 1-\int_0^t g(s)ds\Big]
 \| \nabla u^N(t) \|^2 +(g\circ\nabla u^N)(t)+\| v^N_{t}(t) \|^2 + \| \nabla v^N(t) \|^2 \\
 &+ 2 \alpha \int_{\Omega} v^N(t) u^N(t) dx\Big) \\
&=-\frac{1}{2} g(t)\parallel \nabla u^N(t) \parallel^2 + \frac{1}{2}( g'\circ\nabla u^N )(t).
\end{aligned}\label{e3.6}
\end{equation}
By assumptions (A1)--(A3) we obtain
 \begin{equation}
\begin{aligned}
&\frac{d}{2 dt} \Big( \| u^N_{t} (t) \|^2 + \Big[ 1-\int_0^t g(s)ds\Big]
 \| \nabla u^N(t) \|^2 +(g\circ\nabla u^N)(t)+\| v^N_{t}(t) \|^2
 + \| \nabla v^N (t)\|^2 \\
&+ 2 \alpha \int_{\Omega} v^N(t) u^N(t) dx\Big)\leq 0.
\end{aligned} \label{e3.7}
\end{equation}
Integrating \eqref{e3.7} over $(0, t)$ yields
\begin{equation}
 E_N(t)\leq E_N(0), \label{e3.8}
\end{equation}
where
 \begin{gather*}
 \begin{aligned}
E_N(t)&=\| u^N_{t} (t) \|^2 + \Big[ 1-\int_0^t g(s)ds\Big] \| \nabla u^N(t) \|^2 +(g\circ\nabla u^N)(t)
 +\| v^N_{t}(t) \|^2 + \| \nabla v^N(t) \|^2\\
&\quad + 2 \alpha \int_{\Omega} v^N(t) u^N(t) dx,
\end{aligned}\\
E_N(0)= \| u^N_{1} \|^2 + \| v^N_{1} \|^2 +\| \nabla u^N_0 \|^2
 +\| \nabla v^N_0 \|^2+ 2 \alpha \int_{\Omega} v^N_0 u^N_0 dx.
\end{gather*}
Using Young's and Poincar\'e's inequalities, for $\varepsilon_{1}> 0$, we have
 \begin{equation}
 2\int_{\Omega} v^N(t) u^N(t) dx\geq - \varepsilon_{1} C_p\| \nabla u^N (t) \|^2 - \frac{C_p}{\varepsilon_{1}} \| \nabla v^N (t) \|^2
 \label{e3.9}
 \end{equation}
and
 \begin{equation}
 2 \alpha \int_{\Omega} v^N_0 u^N_0 dx\leq \alpha C_p\| \nabla u^N_0 \|^2 + \alpha C_p \| \nabla v^N_0 \|^2.
 \label{e3.10}
 \end{equation}
Thus
\begin{equation}
 E_N(0)\leq \| u^N_{1} \|^2 + \| v^N_{1} \|^2 +(1+\alpha C_p)\| \nabla u^N_0 \|^2 +(1+\alpha C_p)\| \nabla v^N_0 \|^2.\label{e3.11}
 \end{equation}
Since the sequences $ (u_0^N)_{N \in \mathbb{N}}$,
$ (u_1^N)_{N \in \mathbb{N}}, (v_0^N)_{N \in \mathbb{N}}, (v_1^N)_{N \in \mathbb{N}} $ converge, there exists a positive constant $L_{1}$, independent of $N$, such that
\begin{equation}
E_N(0)\leq L_{1}.\label{e3.12}
\end{equation}
On the other hand,
\begin{equation}
\begin{aligned}
E_N(t)
&\geq \| u^N_{t}(t) \|^2 + \Big[ 1-\int_0^t g(s)ds -\alpha \varepsilon_{1} C_p \Big]
 \| \nabla u^N(t) \|^2 +(g\circ\nabla u^N)(t)+\| v^N_{t} (t) \|^2 \\
 &\quad +\Big( 1-\frac{\alpha C_p}{\varepsilon_{1}}\Big) \| \nabla v^N(t) \|^2.
\end{aligned}\label{e3.13}
\end{equation}
Let us select $\varepsilon_{1}$ such that it satisfies
$\alpha C_p \leq \varepsilon_{1}\leq \frac{l}{\alpha C_p}$.
Then, there exists $C>0$, such that
\begin{equation}
E_N(t)\geq C\left( \| u^N_{t}(t) \|^2 + \| \nabla u^N (t)\|^2
+(g\circ\nabla u^N)(t)+\| v^N_{t}(t) \|^2 + \| \nabla v^N (t)\|^2\right).\label{e3.14}
\end{equation}
From \eqref{e3.8} and \eqref{e3.12}, we deduce
\begin{equation}
\| u^N_{t}(t) \|^2 +\| v^N_{t}(t) \|^2 + \| \nabla u^N(t) \|^2
+ \| \nabla v^N (t)\|^2 +(g\circ\nabla u^N)(t)\leq \frac{L_{1}}{C},\quad
\forall t\leq t_{N}\leq T,
\label{e3.15}
\end{equation}
where $\frac{L_{1}}{C}$ is a constant independent of $t$ and $N$.
\smallskip

\noindent\textbf{A priori estimate 2.}
In \eqref{e3.1}, we substitute $\phi_j $ by $ -\Delta \phi_j$ then
multiply the equation \eqref{e3.1}$_1$ by $a'_j(t)$ and the second equation
\eqref{e3.1}$_2$ by $b'_j(t)$, sum the resulting equations, and use
(A1)--(A3), we arrive at
\begin{align*}
&\frac{d}{2dt}\Big[ \| \nabla u^N_{t}(t) \|^2
+\Big( 1 - \int_0^t g(s)\,ds \Big) \| \Delta u^N(t) \|^2
+ (g \circ \Delta u^N)(t) + \| \nabla v^N_{t}(t) \|^2
+ \| \Delta v^N(t) \|^2 \\
& +2 \alpha \int_{\Omega} \nabla v^N(t) \nabla u^N(t) dx\Big]
-\frac{1}{2} \left( g' \circ \Delta u^N \right)(t)
+\frac{1}{2} g(t) \| \Delta u^N(t) \|^2\leq 0.
\end{align*}

Applying Young's and Poincar\'e's inequalities, we have for $\varepsilon_{3}> 0$,
 \begin{equation*}
 2\int_{\Omega} \nabla v^N(t) \nabla u^N(t) dx
 \geq - \varepsilon_{3} C_p\| \Delta u^N(t) \|^2 - \frac{C_p}{\varepsilon_{3}}\| \Delta v^N(t) \|^2,
 \end{equation*}
and
\[
 2 \int_{\Omega} \nabla v^N_0 \nabla u^N_0 dx
 \leq C_p\| \Delta u^N_0 \|^2 +C_p \| \Delta v^N_0 \|^2.
\]
Then
\begin{equation}
\begin{aligned}
& \| \nabla u^N_{1} \|^2+ \| \nabla v^N_{1} \|^2+\| \Delta u^N_0 \|^2
 + \| \Delta v^N_0 \|^2 +2 \alpha \int_{\Omega} \nabla v^N_0 \nabla u^N_0 dx \\
&\leq \| \nabla u^N_{1} \|^2+ \| \nabla v^N_{1} \|^2
 +(1+\alpha C_p) \| \Delta u^N_0 \|^2
 +(1+\alpha C_p)\| \Delta v^N_0 \|^2.
\end{aligned} \label{e3.16}
\end{equation}
Since the sequences $ (u_0^N)_{N \in \mathbb{N}}$,
$ (u_1^N)_{N \in \mathbb{N}}, (v_0^N)_{N \in \mathbb{N}}, (v_1^N)_{N \in \mathbb{N}} $
converge, there exists a positive constant $L_2$, independent of $N$, such that
\begin{equation}
 \| \nabla u^N_{1} \|^2+ \| \nabla v^N_{1} \|^2
 +\| \Delta u^N_0 \|^2 + \| \Delta v^N_0 \|^2
 +2 \alpha \int_{\Omega} \nabla v^N_0 \nabla u^N_0 dx\leq L_2. \label{e3.17}
\end{equation}
Furthermore,
\begin{align*}
&\| \nabla u^N_{t}(t) \|^2+\Big( 1 - \int_0^t g(s)\,ds \Big) \| \Delta u^N(t) \|^2
 + (g \circ \Delta u^N)(t) + \| \nabla v^N_{t}(t) \|^2+ \| \Delta v^N (t)\|^2 \\
&+2 \alpha \int_{\Omega} \nabla v^N (t)\nabla u^N(t) dx \\
&\geq \| \nabla u^N_{t}(t) \|^2+\Big( 1 - \int_0^t g(s)\,ds
 - \varepsilon_{3} C_p \Big) \| \Delta u^N(t) \|^2 \\
&\quad + (g \circ \Delta u^N)(t) + \| \nabla v^N_{t}(t) \|^2
 +\Big( 1-\frac{\alpha C_p}{\varepsilon_{3}}\Big) \| \Delta v^N(t) \|^2.
\end{align*}

We choose $\varepsilon_{3}$ such that it satisfies
$\alpha C_p \leq \varepsilon_{3}\leq \frac{l}{\alpha C_p}$.
Then there exists $C>0$, such that
\begin{equation}
\begin{aligned}
&\| \nabla u^N_{t}(t) \|^2+\Big( 1 - \int_0^t g(s)\,ds \Big) \| \Delta u^N(t) \|^2
+ (g \circ \Delta u^N)(t) + \| \nabla v^N_{t}(t) \|^2+ \| \Delta v^N(t) \|^2\\
&+2 \alpha \int_{\Omega} \nabla v^N (t)\nabla u^N (t) dx \\
&\geq C \Big( \| \nabla u^N_{t}(t) \|^2+ \| \Delta u^N (t)\|^2
+ (g \circ \Delta u^N)(t) + \| \nabla v^N_{t}(t) \|^2
+ \| \Delta v^N (t)\|^2 \Big).
\end{aligned} \label{e3.18}
\end{equation}
From \eqref{e3.16}, \eqref{e3.17} and \eqref{e3.18}, it follows that
\begin{equation}
\| \nabla u^N_{t}(t) \|^2+ \| \Delta u^N(t) \|^2 + (g \circ \Delta u^N)(t)
+ \| \nabla v^N_{t}(t) \|^2+ \| \Delta v^N(t) \|^2
\leq \frac{L_2}{C}, \; \forall\, t\leq t_{N} \leq T,\label{e3.19}
 \end{equation}
where $\frac{L_2}{C}$ is a constant independent of $t$ and $N$.
\smallskip

\noindent\textbf{A priori estimate 3.}
 By multiplying the first equation in \eqref{e3.1}$_1$ by
 \(a''_j(t)\) and summing the results over \(j\) from 1 to \(k\), we obtain
\begin{align*}
&\int_\Omega |u_{tt}^N(t)|^2 dx \\
& = \int_\Omega \Delta u^N(t) u_{tt}^N(t) \, dx
 - \int_\Omega \int_0^t g(t-s) \Delta u^N(s) u_{tt}^N(t) \, ds\,dx
 - \alpha \int_\Omega v^N(t) u_{tt}^N(t) \, dx\\
& = \int_\Omega \Delta u^N(t) u_{tt}^N(t) \, dx
 + \int_\Omega u_{tt}^N(t) \int_0^t g(t-s) \Big(\Delta u^N(t)- \Delta u^N(s)\Big)\, ds\,dx\\
&\quad -\int_0^t g(s) ds \int_\Omega u_{tt}^N(t) \Delta u^N(t) dx
 - \alpha \int_\Omega v^N(t) u_{tt}^N(t) \, dx\\
&\leq \varepsilon \int_\Omega |u_{tt}^N(t)|^2 dx
 +\frac{1}{4\varepsilon}\int_\Omega |\Delta u^N(t)|^2 \, dx
 + \varepsilon \int_\Omega |u_{tt}^N(t)|^2 dx+\frac{1-l}{4\varepsilon}
 (g \circ \Delta u^N)(t)\\
&\quad +\varepsilon(1-l) \int_\Omega |u_{tt}^N(t)|^2 dx
 + \frac{1-l}{4\varepsilon} \int_\Omega |\Delta u^N(t)|^2dx
 + \varepsilon \int_\Omega |u_{tt}^N(t)|^2 dx
 +\frac{\alpha^2}{4\varepsilon}\int_\Omega |v^N(t)|^2 dx \\
&\leq \varepsilon (4-l)\int_\Omega |u_{tt}^N(t)|^2 dx
 +\frac{2-l}{4\varepsilon} \int_\Omega |\Delta u^N(t)|^2dx+
 \frac{1-l}{4\varepsilon} (g \circ \Delta u^N)(t)
 +\frac{\alpha^2}{4\varepsilon}\int_\Omega |v^N(t)|^2 dx.
\end{align*}
Hence, for some $c >0$, we have
\begin{align}
\left(1-\varepsilon(4-l)\right) \int_\Omega |u_{tt}^N(t)|^2 dx
\leq \frac{c}{4\varepsilon}(L_1+L_2).
\end{align}
We take $\varepsilon >0$ small enough, so
\begin{align}
\|u_{tt}^N(t)\|^2 \leq \tilde{c}(L_1+L_2)=L_3. \label{e3.20}
\end{align}
Similarly, we obtain $\|v_{tt}^N(t)\|^2 \leq L_4$, where $L_{3}, L_{4} $
are constants independent of $t$ and $N$.

From \eqref{e3.15}, \eqref{e3.19} and \eqref{e3.20}, we conclude that
\begin{equation} \label{e3.21}
\begin{gathered}
(u^{N})_{N} \text{ and } (v^N)_{N} \text{ are bounded in }
L^\infty((0,T), H^2(\Omega) \cap H_0^1 (\Omega)),\\
(u^N_{t})_{N} \text{ and } (v^N_{t})_{N} \text{ are bounded in }
L^\infty((0,T), H^1_0(\Omega)), \\
(u^N_{tt})_{N} \text{ and } (v^N_{tt})_{N} \text{ are bounded in }
 L^\infty((0,T), L^2(\Omega)).
\end{gathered}
\end{equation}

\noindent\textbf{Step 3.}
From \eqref{e3.21}, there exist subsequences of $(u^{N})_{N}$ and $(v^{N})_{N}$
(still denoted by $(u^{N})_{N}$ and $(v^{N})_{N})$, and two functions
$u, v : \Omega \times [0, T) \to \mathbb{R}$, such that
\begin{equation}
\begin{gathered}
u^{N} \overset{\star}{\rightharpoonup} u \text{ and }
 v^{N} \overset{\star}{\rightharpoonup} v \quad \text{in }
 L^\infty((0,T), H^2(\Omega) \cap H_0^1 (\Omega)), \\
u_{t}^{N} \overset{\star}{\rightharpoonup} u_{t} \text{ and }
 v_{t}^{N} \overset{\star}{\rightharpoonup} v_{t} \text{ in } L^\infty((0,T),
 H^1_0(\Omega)), \\
 u_{tt}^{N} \overset{\star}{\rightharpoonup} u_{tt} \text{ and }
 v_{tt}^{N} \overset{\star}{\rightharpoonup} v_{tt} \text{ in }
 L^\infty((0,T), L^2(\Omega)).
\end{gathered}\label{e3.22}
\end{equation}
Thanks to convergences \eqref{e3.2} and \eqref{e3.22}, taking $N \to +\infty$
in \eqref{e3.1}, we obtain
\begin{equation}
\begin{gathered}
\int_\Omega u_{tt}\phi_j dx + \int_\Omega \nabla u .\nabla \phi_j dx
- \int_\Omega \int_0^t g(t-s) \nabla u(s) .\nabla \phi_j\, ds\, dx
+ \alpha \int_\Omega v \phi_j \, dx = 0,\\
\int_\Omega v_{tt} \phi_j dx + \int_\Omega \nabla v.\nabla \phi_j dx
+ \alpha \int_\Omega u \phi_j \, dx = 0,\quad 1\leq j\leq N.
\end{gathered} \label{e3.23}
\end{equation}
 Consequently, for every $ \phi \in H_0^1(\Omega)\cap H^2(\Omega)$,
\begin{gather*}
\int_\Omega \Big[ u_{tt}- \Delta u + \int_0^t g(t-s) \Delta u(s)\, ds\,dx
 + \alpha v\Big] \phi \, dx = 0,\\
\int_\Omega [ v_{tt} - \Delta v + \alpha u] \phi \, dx = 0.
\end{gather*}
Therefore, $(u,v)$ solves the equations in \eqref{e1.1}.
\smallskip

\noindent\textbf{Step 4. Initial conditions.}
First, by Lions' lemma \cite{Lions} and since
\begin{gather*}
u^{N} \overset{\star}{\rightharpoonup} u \quad \text{in }
 L^\infty((0,T), H^2(\Omega) \cap H_0^1 (\Omega)), \\
u_{t}^{N} \overset{\star}{\rightharpoonup} u_{t} \quad \text{in }
 L^\infty((0,T), H^1_0(\Omega)),
\end{gather*}
we deduce, up to a subsequence, that
\begin{equation}
u^{N} \to u \quad \text{in } C([0, T], H^1_0(\Omega)).\label{e3.24}
\end{equation}
Therefore, \( u^{N}(\cdot, 0) \) is defined and
\[
u^{N}(\cdot, 0) \to u(\cdot, 0) \quad \text{in } H^1_0(\Omega).
\]
But,
\[ u^{N}(\cdot, 0) = u^{N}_0 \to u_0 \,\, \text{in}\,\, H^2(\Omega) \cap H^1_0(\Omega).\] Then \( u(\cdot, 0) = u_0 \, \). Similarly, we obtain $v(\cdot, 0)=v_0. $

Noting that
\[
 u_{tt}^{N} \overset{\star}{\rightharpoonup} u_{tt} \quad \text{and} \quad
 v_{tt}^{N} \overset{\star}{\rightharpoonup} v_{tt} \quad \text{in }
 L^\infty((0,T), L^2(\Omega)),
\]
we obtain $ u_t , v_{t} \in C([0, T), L^2(\Omega))$, as in lions \cite{Lions},
and thus \( u_t(\cdot, 0) \) has a meaning. Furthermore
\[
u^N_t(\cdot, 0) \to u_t(\cdot, 0) \quad \text{and} \quad
v^N_t(\cdot, 0) \to v_t(\cdot, 0) \quad\text{in } L^2(\Omega).
\]
But,
\[ u^N_t(\cdot, 0) = u^N_1 \to u_1 \quad \text{and} \quad
 v^N_t(\cdot, 0) = v^N_1 \to v_1 \quad \text{in } H^1_0(\Omega).
\]
Hence,
$u_t(x, 0) = u_1(x)$ and $v_t(x, 0) = v_1(x)$.
\smallskip

\noindent\textbf{Uniqueness.}
Assume that problem \eqref{e1.1} has two weak solutions, \( (u, v) \) and
\( (\overline{u},\overline{v}) \). Then,
\( (\mathbf{u}, \mathbf{v}) = (u - \overline{u}, v - \overline{v}) \) satisfies
\begin{gather}
 \mathbf{u}_{tt} - \Delta \mathbf{u} + \int_0^t g(t-s) \Delta \mathbf{u} \, ds
 + \alpha \mathbf{v} = 0, \quad \text{in } \Omega \times (0, T), \label{e3.25}\\
 \mathbf{v}_{tt} - \Delta \mathbf{v} + \alpha \mathbf{u} = 0, \quad
 \text{in } \Omega \times (0, T), \label{e3.26} \\
 \mathbf{u}(x,t) = \mathbf{v}(x,t) = 0, \quad \text{on } \partial\Omega \times (0, T),
 \label{e3.27}\\
 \mathbf{u}(x,0) = \mathbf{u}_t(x,0) =0,\, \mathbf{v}(x,0)=\mathbf{v}_t(x,0) = 0,
 \quad \text{in } \Omega. \label{e3.28}
\end{gather}
Multiplying \eqref{e3.25} by $\mathbf{u}_{t}$ and \eqref{e3.26} by $\mathbf{v}_{t}$,
integrating over $\Omega$, and adding the two resulting equations, we find
\begin{equation}
\begin{aligned}
&\frac{1}{2}\frac{d}{dt} \Big( \| \mathbf{u}_{t}(t) \| ^2
 + \Big[ 1-\int_0^t g(s) ds \Big] \| \nabla \mathbf{u} (t)\|^2
 +(g\circ \nabla \mathbf{u})(t) \\
&+2 \alpha \int_\Omega \mathbf{v}\textbf{ u}\,dx
 +\| \mathbf{v}_{t}(t) \|^2 + \| \nabla \mathbf{v} (t)\|^2 \Big)\\
& = -\frac{1}{2} g(t)\parallel \nabla \mathbf{u}(t) \parallel^2
 + \frac{1}{2}( g'\circ\nabla \mathbf{u}) (t).
\end{aligned} \label{e3.29}
\end{equation}
Using assumptions (A1)--(A3), and integrating over $(0, t),\, t\leq T$, we obtain
\[
 \| \mathbf{u}_{t}(t) \| ^2+ \Big[ 1-\int_0^t g(s) ds \Big]
 \| \nabla \mathbf{u}(t) \|^2 +(g\circ \nabla \mathbf{u})(t)
 +2 \alpha \int_\Omega \mathbf{v}\textbf{ u}dx
 +\| \mathbf{v}_{t} (t)\|^2 + \| \nabla \mathbf{v}(t) \|^2 \leq 0.
\]
 From Young and Poincar\'e's inequalities, we obtain
\begin{align*}
&\| \mathbf{u}_{t}(t) \| ^2+ \Big[ 1-\int_0^t g(s) ds
 -\alpha \varepsilon_{1} C_p \Big] \| \nabla \mathbf{u}(t) \|^2
 +(g\circ \nabla \mathbf{u})(t)\\
&+\| \mathbf{v}_{t}(t) \|^2 + \Big( 1-\frac{\alpha C_{p}}{\varepsilon_{1} }\Big)
\| \nabla \mathbf{v}(t) \|^2 \leq 0.
\end{align*}
We select $\varepsilon_{1}$ such that it satisfies\,
$\alpha C_p \leq \varepsilon_{1}\leq \frac{l}{\alpha C_p}$. Then
\[
 \| \mathbf{u}_{t} \|^2 + \| \nabla \mathbf{u} \|^2
 +(g\circ\nabla \mathbf{u})(t)+\| \mathbf{v}_{t} \|^2 + \| \nabla \mathbf{v} \|^2=0.
\]
Hence, \( \mathbf{u}_t(\cdot, t) = \mathbf{v}_t(\cdot, t) = 0 \) and
\( \nabla \mathbf{u}(\cdot, t) = \nabla\mathbf{v}(\cdot, t) = 0 \), for all
\( t \in (0, T) \). This implies that \( \textbf{ u} = \textbf{ v} = 0 \) on
\( \Omega \times (0, T) \), given that \( \mathbf{u} = \mathbf{v} = 0 \) on
\( \partial\Omega \times (0, T) \). This establishes the uniqueness.
\end{proof}

\begin{thebibliography}{99}
\bibitem[1]{Assila1} M. Aassila; \emph{A note on the boundary stabilization of a compactly coupled system of wave equations}. Applied Mathematics Letters 12.3 (1999): 19-24.
\bibitem[2]{Assila2} Mohammed Aassila; \emph{Strong asymptotic stability of a compactly coupled system of wave equations}. Applied Mathematics Letters 14.3 (2001): 285-290.
\bibitem[3]{AkilM} Mohammad, Akil, et al.; \emph{Stability results of coupled wave models with locally memory in a past history framework via nonsmooth coefficients on the interface}. Mathematical Methods in the Applied Sciences 44.8 (2021): 6950-6981.
\bibitem[4]{Beniani} Abderrahmane Beniani,  Noureddine Taouaf,  Abbes Benaissa; \emph{Well-posedness and exponential stability for coupled Lam\'e system with viscoelastic term and strong damping}. Computers Mathematics with Applications 75.12 (2018): 4397-4404.
\bibitem[5]{Berrimi2004} Said Berrimi, Salim A. Messaoudi; \emph{Exponential decay of solutions to a viscoelastic equation with nonlinear localized damping}. Electronic Journal of Differential Equations, 2004 (2004), No. 88, 1=10.
\bibitem[6]{Berrimi2006} Said Berrimi, Salim A. Messaoudi; \emph{Existence and decay of solutions of a viscoelastic equation with a nonlinear source}. Nonlinear Analysis: Theory, Methods Applications 64.10 (2006): 2314-2331.
\bibitem[7]{CavalcantiDS} M. M. Cavalcanti, VN Domingos Cavalcanti, J. Ferreira; \emph{Existence and uniform decay for a non-linear viscoelastic equation with strong damping}. Mathematical methods in the applied sciences 24.14 (2001): 1043-1053.
\bibitem[8]{CavalcantiFerreira} M. M. Cavalcanti, VN Domingos Cavalcanti, J. Ferreira; \emph{Exponential decay for the solution of semilinear viscoelastic wave equation with localized damping}. Electronic J. Diff. Eqns., 2002 (2002) No. 44, 1-14.
\bibitem[9]{Cavalcanti2003} M. M. Cavalcanti, Higidio Portillo Oquendo; \emph{Frictional versus viscoelastic damping in a semilinear wave equation}. SIAM journal on control and optimization 42.4 (2003): 1310-1324.
\bibitem[10]{Dafermos1} Constantine M. Dafermos; \emph{An abstract Volterra equation with applications to linear viscoelasticity}. Journal of Differential Equations 7.3 (1970): 554-569.
\bibitem[11]{Dafermos2} Constantine M. Dafermos; \emph{Asymptotic stability in viscoelasticity}. Archive for rational mechanics and analysis 37 (1970): 297-308.
\bibitem[12]{FengBaoweiHaiyan} Baowei Feng, Haiyan Li; \emph{Decay rates for a coupled viscoelastic Lam\'e system with strong damping}. Mathematical Modelling and Analysis 25.2 (2020): 226-240.
\bibitem[13]{GASM} Aissa Guesmia, Salim A. Messaoudi, Mostafa Zahri; \emph{General decay of solutions of a weakly coupled abstract evolution equations with one finite memory control}. Journal of Applied Analysis, Computation 14.6 (2024): 3539-3557.
\bibitem[14]{Hao} Jianghao Hao,  Fei Wang; \emph{General decay rate for weak viscoelastic wave equation with Balakrishnan-Taylor damping and time-varying delay}. Computers Mathematics with Applications 78.8 (2019): 2632-2640.
\bibitem[15]{JKJT} Kun-Peng Jin, Jin Liang,  Ti-Jun Xiao; \emph{Coupled second order evolution equations with fading memory: Optimal energy decay rate}. Journal of Differential Equations 257.5 (2014): 1501-1528.
\bibitem[16]{KVBR} Vilmos Komornik, Bopeng Rao; \emph{Boundary stabilization of compactly coupled wave equations}. Asymptotic Analysis 14.4 (1997): 339-359.
\bibitem[17]{Li2011} Fushan Li, Zengqin Zhao, Yanfu Chen; \emph{Global existence uniqueness and decay estimates for nonlinear viscoelastic wave equation with boundary dissipation}. Nonlinear Analysis: Real World Applications 12.3 (2011): 1759-1773.
\bibitem[18]{Lions} J. L. Lions; \emph{Quelques methodes de resolution des problemes aux limites non lineaires}. Dunford/Gauthier-Villars (1969).
\bibitem[19]{LRF} Renato F.C. Lobato, et al.; \emph{Optimal polynomial decay to coupled wave equations and its numerical properties} Journal of Applied Mathematics 2014.1 (2014): 897080.
\bibitem[20]{SMessaoudi1} Salim A. Messaoudi; \emph{General decay of solutions of a viscoelastic equation}. Journal of Mathematical Analysis and Applications 341.2 (2008): 1457-1467.
\bibitem[21]{SMessaoudi2} Salim Messaoudi; \emph{General decay of the solution energy in a viscoelastic equation with a nonlinear source}. Nonlinear Analysis: Theory, Methods Applications 69.8 (2008): 2589-2598.
\bibitem[23]{MessaoudiMustafa} Salim A. Messaoudi,  Muhammad Islam Mustafa; \emph{On the control of solutions of viscoelastic equations with boundary feedback}. Nonlinear Analysis: Real World Applications 10.5 (2009): 3132-3140.
\bibitem[24]{MessaoudiWaled} Salim A. Messaoudi, Waled Al-Khulaifi; \emph{General and optimal decay for a quasilinear viscoelastic equation}. Applied Mathematics Letters 66 (2017): 16-22.
\bibitem[25]{MTatar1} Salim A. Messaoudi,  N. E. Tatar; \emph{Global existence and asymptotic behavior for nonlinear viscoelastic problem}. Mathematical Sciences Research Journal 7 (2003): 136-149.
\bibitem[26]{MTatar2} Salim A. Messaoudi,  Nasser-eddine Tatar; \emph{Global existence and uniform stability of solutions for a quasilinear viscoelastic problem}. Mathematical Methods in the Applied Sciences 30 (6) (2007).
\bibitem[27]{MM181} Muhammad I. Mustafa; \emph{Optimal decay rates for the viscoelastic wave equation}. Mathematical Methods in the Applied Sciences 41.1 (2018): 192-204.
\bibitem[28]{MM} Muhammad I. Mustafa, Salim A. Messaoudi; \emph{General stability result for viscoelastic wave equations}. Journal of Mathematical Physics 53.5 (2012).
\bibitem[30]{Munoz} Jaime E. Munoz Rivera, Doherty Andrade; \emph{Exponential decay of non-linear wave equation with a viscoelastic boundary condition}. Mathematical Methods in the Applied Sciences 23.1 (2000): 41-61.
\bibitem[31]{ORHC} Rafael Lima Oliveira, Higidio Portillo Oquendo, Celene Buriol; \emph{Exact decay rates for weakly coupled systems with indirect damping by fractional memory effects}. Mathematische Nachrichten 296.10 (2023): 4610-4633.
\bibitem[32]{PJS} Jong Yeoul Park, Sun Hye Park; \emph{General decay for quasilinear viscoelastic equations with nonlinear weak damping}. Journal of Mathematical Physics 50.8 (2009).
\bibitem[33]{SaidH2011} Belkacem Said-Houari, Salim A. Messaoudi, Aissa Guesmia; \emph{General decay of solutions of a nonlinear system of viscoelastic wave equations}. Nonlinear Differential Equations and Applications NoDEA 18.6 (2011): 659-684.
\bibitem[34]{WSH} Shun-Tang Wu; \emph{General decay for a wave equation of Kirchhoff type with a boundary control of memory type}. Boundary Value Problems (2011): 1-15.
\end{thebibliography}
\end{document}
